\chapter{Loop Analysis}
\label{ch:ch18-loop-analysis}

\section{Seeing Loops in Graphs}
\label{sec:ch18-seeing-loops}

Optimizers spend most of their effort on loops because programs spend most of their time in loops.
Motion, unrolling, tiling, vectorization, and parallelization all begin by answering the same preliminary questions.
Where does each loop start, which blocks belong to it, how do loops nest, and what changes from one iteration to the next.
Guessing from source syntax is unreliable, because lowering turns structured loops into flat branches and irreducible control flow can defeat syntax entirely.
\term{Loop analysis} answers these questions on the control-flow graph itself.
It defines a loop as a graph shape, finds every instance of that shape, organizes the results into nesting, and characterizes how values evolve across iterations.
The definitions must be exact, because every later transformation trusts them.
A block wrongly assigned to a loop gets its code moved to the wrong place.
A missed back edge hides an entire loop from optimization.
This chapter builds the standard pipeline: dominators identify back edges, back edges identify natural loops, nesting organizes them, and induction analysis describes iteration.
It uses the loop facts in Chapter~\ref{ch:ch19-loop-transformations} and supplies the regions that Chapter~\ref{ch:ch17-code-motion} requires.

\section{Back Edges and Natural Loops}
\label{sec:ch18-back-edges}

A \term{back edge} is an edge $u \to h$ whose target $h$ dominates its source $u$.
Domination means every path from the entry to $u$ passes through $h$, so control cannot reach $u$ without first passing $h$.
Returning to $h$ therefore repeats a region that starts at $h$, which is exactly the intuition a loop should capture.
The \term{loop header} is the target $h$, and the source $u$ is often called the latch.
Given a back edge $u \to h$, the \term{natural loop} consists of $h$ plus every block that can reach $u$ without passing through $h$.
That definition is constructive: start from $u$, walk predecessors, stop at $h$, and collect everything visited.
The header dominates every block in its loop, which follows from the back-edge property and makes the header the single entry of the loop.
Single entry matters enormously.
It provides a single entry at which a preheader can be inserted.
A preheader can execute even when a header test skips the body, so hoisting still requires a speculation-safety or must-execute argument.
Some graphs contain loops with several entries, usually produced by unstructured jumps.
Those \term{irreducible} loops have no single header, and most loop optimizations decline to handle them.
Compilers either leave such regions alone or restructure them first, accepting that rare shapes get slower code.
Listing~\ref{lst:ch18-nested} shows the running example for this chapter: two nested counted loops after lowering.

\begin{irlisting}[caption={Two nested loops with headers at blocks 2 and 3 and back edges from 4 and 5.},label={lst:ch18-nested}]
1: goto 2
2: i = phi(0, inext) ; outer header
   if i < 10 goto 3 else goto 6
3: j = phi(0, jnext) ; inner header
   if j < 20 goto 4 else goto 5
4: body(i, j)
   jnext = j + 1
   goto 3
5: inext = i + 1
   goto 2
6: ret
\end{irlisting}

The inner back edge runs from block $4$ to block $3$, and the outer back edge runs from block $5$ to block $2$.
The inner header differs from the outer header, and the inner loop body sits strictly inside the outer one.
Algorithm~\ref{alg:ch18-natural-loops} turns these observations into a procedure.

\begin{algorithm}
  \caption{Natural loop discovery from back edges}\label{alg:ch18-natural-loops}
  \begin{pseudocode}
    \Require entry-reachable CFG, dominator sets $\dom(B)$ for every block $B$
    \Ensure the set of natural loops, each as a header plus member blocks
    \State $Loops \gets$ empty set
    \For{each edge $u \to h$ in the CFG}
      \If{$h \in \dom(u)$}
        \State $Body \gets \{h, u\}$
        \State $Work \gets \{u\}\setminus\{h\}$
        \While{$Work \neq \emptyset$}
          \State remove some $x$ from $Work$
          \For{each predecessor $p \in \Pred(x)$}
            \If{$p \notin Body$}
              \State add $p$ to $Body$ and to $Work$
            \EndIf
          \EndFor
        \EndWhile
        \State add $(h, Body)$ to $Loops$
      \EndIf
    \EndFor
  \end{pseudocode}
\end{algorithm}

The algorithm assumes dominator sets are already computed, since testing $h \in \dom(u)$ is the heart of back-edge detection.
The predecessor walk never passes through the header, because the header is already in $Body$ and visited blocks are not re-added.
Each walk therefore collects exactly the blocks that reach the latch without leaving through the header.

\section{Worked Example: Dominators to Loops}
\label{sec:ch18-worked-loops}

We run Algorithm~\ref{alg:ch18-natural-loops} on Listing~\ref{lst:ch18-nested} by hand.
First we compute dominator sets with the entry dominating itself and every other block dominated by the intersection of its predecessors plus itself.
Only block $1$ dominates the entry, so $\dom(1) = \{1\}$; the entry itself dominates every reachable block.
Block $2$ is reachable only through block $1$, so $\dom(2) = \{1, 2\}$.
Block $3$ is reachable only through block $2$, so $\dom(3) = \{1, 2, 3\}$.
Block $4$ is reachable only through block $3$, so $\dom(4) = \{1, 2, 3, 4\}$.
Block $5$ is reachable only through block $3$, so $\dom(5) = \{1, 2, 3, 5\}$.
Block $6$ is reachable only through block $2$, so $\dom(6) = \{1, 2, 6\}$.
We check every edge for the domination property.
The edge $4 \to 3$ qualifies because $3 \in \dom(4)$, and the edge $5 \to 2$ qualifies because $2 \in \dom(5)$.
No other edge targets a dominator of its source: edge $3 \to 5$ fails because $5 \notin \dom(3)$, and edge $2 \to 3$ fails because $3 \notin \dom(2)$.
For the inner back edge, the walk starts with $Body = \{3, 4\}$ and $Work = \{4\}$.
Block $4$ has the single predecessor $3$, which is already in $Body$, so the walk ends with the inner loop $\{3, 4\}$.
For the outer back edge, the walk starts with $Body = \{2, 5\}$ and $Work = \{5\}$.
Block $5$ has predecessor $3$, which joins $Body$ and $Work$.
Block $3$ has predecessors $2$ and $4$, of which $2$ is already present and $4$ joins both sets.
Block $4$ has predecessor $3$, already present, so the walk ends with the outer loop $\{2, 3, 4, 5\}$.
Table~\ref{tab:ch18-dom-loops} summarizes the dominators, the back-edge test, and the resulting membership.
Nesting follows by inclusion: the inner loop is strictly contained in the outer one, so the \term{loop nesting tree} has the outer loop as parent.
Figure~\ref{fig:ch18-nested} draws the graph so the nesting is visible.

\begin{table}
  \centering
  \caption{Dominator sets, back-edge tests, and loop membership for Listing~\ref{lst:ch18-nested}.}
  \label{tab:ch18-dom-loops}
  \begin{tabular}{llll}
    \toprule
    Block & Dominators & Back edge & Loop membership \\
    \midrule
    1 & 1 & none & none \\
    2 & 1, 2 & target of 5 to 2 & outer header \\
    3 & 1, 2, 3 & target of 4 to 3 & outer, inner header \\
    4 & 1, 2, 3, 4 & source of 4 to 3 & outer, inner \\
    5 & 1, 2, 3, 5 & source of 5 to 2 & outer \\
    6 & 1, 2, 6 & none & none \\
    \bottomrule
  \end{tabular}
\end{table}

\begin{figure}
  \centering
  \begin{tikzpicture}[node distance=12mm and 20mm]
    \node[basic block] (ch18n1) {1: goto 2};
    \node[basic block,below=of ch18n1] (ch18n2) {2: outer header};
    \node[basic block,below left=of ch18n2] (ch18n3) {3: inner header};
    \node[basic block,below=of ch18n3] (ch18n4) {4: body};
    \node[basic block,below right=of ch18n2] (ch18n5) {5: outer latch};
    \node[basic block,below=of ch18n5] (ch18n6) {6: ret};
    \draw[cfg edge] (ch18n1) -- (ch18n2);
    \draw[cfg edge] (ch18n2) -- (ch18n3);
    \draw[cfg edge] (ch18n2) -- (ch18n6);
    \draw[cfg edge] (ch18n3) -- (ch18n4);
    \draw[cfg edge] (ch18n3) -- (ch18n5);
    \draw[cfg edge] (ch18n4) -- (ch18n3);
    \draw[back edge] (ch18n5) -- (ch18n2);
  \end{tikzpicture}
  \caption{The inner loop sits inside the outer loop, and each back edge returns to the header that dominates its source.}
  \label{fig:ch18-nested}
\end{figure}

\section{Trip Counts and Induction Variables}
\label{sec:ch18-induction}

Knowing where a loop is matters less than knowing how it behaves.
A \term{trip count} says how many times the loop iterates, and an \term{induction variable} is a variable that steps by a fixed amount each iteration.
The outer loop of Listing~\ref{lst:ch18-nested} counts \code{i} from $0$ to $9$ in steps of $1$, so its trip count is $10$.
The inner loop counts \code{j} from $0$ to $19$, so its trip count is $20$ and the body runs $200$ times in total.
Both counts follow from the bounds and steps visible in the IR, and both are exact because the bounds are constants and the steps never change.
A \term{basic induction variable} is defined once per iteration by adding or subtracting a loop-invariant amount, like \code{inext = i + 1}.
A \term{derived induction variable} is an affine function of a basic one, like \code{u = i * 4} from Chapter~\ref{ch:ch17-code-motion}.
Derived variables matter because address arithmetic is usually derived from a counter, and recognizing the derivation lets later passes replace multiplication with addition.
The analysis walks SSA definitions inside the loop and classifies each one.
Definitions from outside constants are invariant, single-step updates of loop-carried values are basic, and affine combinations of classified values are derived.
Anything else is simply variant, and the analysis says so rather than guessing.
Derived variables deserve one more observation.
The expression \code{u = i * 4} multiplies a basic induction variable by a constant, which strength reduction later rewrites into an addition stepped by four each iteration.
That rewrite is only valid because the analysis proved the affine form first.
Without the proof, replacing a multiplication by an accumulation changes overflow behavior on targets where the two differ.
Trip counts need the same honesty.
A data-dependent bound can admit a symbolic exact count, while an early exit or possible counter overflow may invalidate the simple counted-loop formula.
When the available reasoning cannot establish a count, report \enquote{unknown}; lack of a constant is not proof that no exact symbolic count exists.
An unknown static count can still permit unrolling or tiling with runtime bounds and a cleanup loop.
Motion and redundancy removal need their own legality facts rather than an exact trip count.

\section{Implementation Notes}
\label{sec:ch18-implementation}

The set-based dominator iteration costs $O(R(N+E)W)$ for $R$ rounds, $N$ blocks, $E$ edges, and cost $W$ per set operation.
Reverse-postorder traversal can reduce the number of rounds but does not remove the set-operation cost.
% TODO(cite): primary reference for near-linear dominator computation.
Each back edge triggers one predecessor walk bounded by $O(E)$, so loop discovery over $b$ back edges costs $O(b \cdot E)$ in the worst case and far less in practice.
Sorting $L$ discovered loops costs $O(L\log L)$, but constructing a nesting forest also requires membership or containment queries.
Merge natural loops with a common header under a documented convention before assigning parent loops; sorting alone does not establish containment.
Induction analysis is linear in the number of SSA definitions inside loops.
Production compilers reuse one dominator computation across many clients rather than recomputing it per analysis.
They also maintain the loop tree incrementally across transformations, repairing only the affected subtrees instead of rediscovering every loop.
Irreducible flow needs a policy, not just an algorithm.
The common choice marks irreducible regions, skips loop transformations there, and keeps the code correct at the price of missing rare optimizations.
Guard this choice with a counter in the compiler's statistics so that an unexpected spike of irreducible loops points back to the frontend pass that created them.

\section{Exercises}
\label{sec:ch18-exercises}

\begin{enumerate}
  \item \label{ex:ch18-find-loops}
  Add the edge $5 \to 3$ to Listing~\ref{lst:ch18-nested} alongside the existing $5 \to 2$.
  Recompute the dominator sets, test every edge with the back-edge condition, and run Algorithm~\ref{alg:ch18-natural-loops} by hand.
  State the resulting loop set and nesting, and explain which edge, if any, fails the domination test.
  \item \label{ex:ch18-irreducible}
  Draw a CFG with entry $A$, blocks $B$ and $C$ each branching to each other, entries $A \to B$ and $A \to C$, and a shared exit.
  Compute dominator sets and show that neither $B \to C$ nor $C \to B$ is a back edge.
  Explain why no natural loop covers the cycle, and state what an optimizer that only handles natural loops should do with this region.
  \item \label{ex:ch18-classify-iv}
  Given a loop with header definitions \code{i = phi(0, i2)}, \code{k = phi(0, k2)}, \code{i2 = i + 2}, \code{k2 = k + i2}, \code{m = k * 3}, and \code{n = load p}, classify each of \code{i}, \code{k}, \code{m}, and \code{n} as invariant, basic induction, derived induction, or variant.
  Justify each classification from the definitions in Section~\ref{sec:ch18-induction}.
  Give the trip count when the exit tests \code{i2 < 100}, showing the arithmetic.
\end{enumerate}

\section{Further reading}
\label{sec:ch18-further-reading}

Morgan's \emph{Building an Optimizing Compiler} develops control-flow analysis including loops and their use in optimization~\cite{morgan-1998-optimizing}.
Read it for how loop structure connects to the transformations that consume it.
Appel's \emph{Modern Compiler Implementation in ML} covers the graph algorithms behind this chapter in implementable form~\cite{appel-1998-modern}.
Use it when turning the pseudocode here into working code.
Safonov's \emph{Trustworthy Compilers} discusses the obligations that transformations over these structures must meet~\cite{safonov-2010-trustworthy}.
Consult it for why a wrong loop membership is a correctness bug rather than a missed opportunity.
